% ======================================================
%  Sekcja 4 --- Ci\k{a}gi i Wzorce
%  8 zada\'{n} | Arytmetyczne, Geometryczne, Fibonacciego, Literowe
% ======================================================

% -- 26 --
\ExerciseSlide[\CategoryBadge[SeqColor!20]{Arytmetyczny}]
  {1}{30 s}{$2,\ 4,\ 6,\ 8,\ \mathbf{?}$}

\AnswerSlide[\CategoryBadge[SeqColor!20]{Arytmetyczny}]{10}{
  {\footnotesize
    Ci\k{a}g arytmetyczny, r\'{o}\. znica $d = +2$.
    \[ 8 + 2 = \mathbf{10} \]
  }
}

% -- 27 --
\ExerciseSlide[\CategoryBadge[SeqColor!20]{Fibonacciego}]
  {1}{30 s}{$1,\ 1,\ 2,\ 3,\ 5,\ 8,\ \mathbf{?}$}

\AnswerSlide[\CategoryBadge[SeqColor!20]{Fibonacciego}]{13}{
  {\footnotesize
    Ci\k{a}g Fibonacciego: ka\. zdy wyraz = suma dw\'{o}ch poprzednich.
    \[ 5 + 8 = \mathbf{13} \]
  }
}

% -- 28 --
\ExerciseSlide[\CategoryBadge[SeqColor!20]{Geometryczny}]
  {2}{60 s}{$2,\ 6,\ 18,\ 54,\ \mathbf{?}$}

\AnswerSlide[\CategoryBadge[SeqColor!20]{Geometryczny}]{162}{
  {\footnotesize
    Ci\k{a}g geometryczny, iloraz $q = 3$.
    \[ 54 \times 3 = \mathbf{162} \]
  }
}

% -- 29 --
\ExerciseSlide[\CategoryBadge[SeqColor!20]{Kwadraty}]
  {2}{60 s}{$1,\ 4,\ 9,\ 16,\ 25,\ \mathbf{?}$}

\AnswerSlide[\CategoryBadge[SeqColor!20]{Kwadraty}]{36}{
  {\footnotesize
    Kolejne kwadraty liczb naturalnych: $1^2, 2^2, 3^2, 4^2, 5^2, \ldots$
    \[ 6^2 = \mathbf{36} \]
  }
}

% -- 30 --
\ExerciseSlide[\CategoryBadge[SeqColor!20]{R\'{o}\. znice}]
  {2}{90 s}{$3,\ 7,\ 13,\ 21,\ 31,\ \mathbf{?}$}

\AnswerSlide[\CategoryBadge[SeqColor!20]{R\'{o}\. znice}]{43}{
  {\footnotesize
    R\'{o}\. znice: $4, 6, 8, 10, 12, \ldots$ (rosn\k{a}ce parzyste).
    \[ 31 + 12 = \mathbf{43} \]
  }
}

% -- 31 --
\ExerciseSlide[\CategoryBadge[SeqColor!20]{R\'{o}\. znice}]
  {3}{2 min}{$1,\ 2,\ 4,\ 7,\ 11,\ 16,\ \mathbf{?}$}

\AnswerSlide[\CategoryBadge[SeqColor!20]{R\'{o}\. znice}]{22}{
  {\footnotesize
    R\'{o}\. znice: $1, 2, 3, 4, 5, 6, \ldots$ (kolejne liczby naturalne).
    \[ 16 + 6 = \mathbf{22} \]
  }
}

% -- 32 --
\ExerciseSlide[\CategoryBadge[SeqColor!20]{Pierwsze}]
  {3}{2 min}{$2,\ 3,\ 5,\ 7,\ 11,\ 13,\ \mathbf{?}$}

\AnswerSlide[\CategoryBadge[SeqColor!20]{Pierwsze}]{17}{
  {\footnotesize
    Kolejne liczby pierwsze.
    \[ \text{Po 13 nast\k{e}pna liczba pierwsza to } \mathbf{17} \]
    (14 = 2\texttimes7, 15 = 3\texttimes5, 16 = $2^4$ --- \. zadna nie jest pierwsza)
  }
}

% -- 33 --
\ExerciseSlide[\CategoryBadge[SeqColor!20]{Literowy}]
  {3}{3 min}{$\mathrm{A,\ B,\ D,\ G,\ K,\ \mathbf{?}}$}

\AnswerSlide[\CategoryBadge[SeqColor!20]{Literowy}]{P}{
  {\footnotesize
    Pozycje w alfabecie: $1, 2, 4, 7, 11, \ldots$\par
    R\'{o}\. znice mi\k{e}dzy wyrazami: $+1, +2, +3, +4, +5, \ldots$
    \[ 11 + 5 = 16 \Rightarrow
       \text{16.~litera alfabetu} = \mathbf{P} \]
  }
}
